\item Un alambre de densidad $\rho$ se afila de modo que su área de sección transversal varía con $x$ de acuerdo con
$$ A = 1.00 \times 10^{-5} x + 1.00 \times 10^{-6} $$
donde $A$ está en metros cuadrados y $x$ en metros. La tensión en el alambre es $T$.
\begin{enumerate}
    \item Deduzca una relación para la rapidez de una onda como función de la posición.
    \item ¿Qué pasaría si? Suponga que el alambre es de aluminio y está bajo una tensión $T = 24.0\text{ N}$. Determine la rapidez de onda en el origen y en $x = 10.0\text{ m}$.
\end{enumerate}